%% Cover letter accompanying the submission of
%% "Bayesian fusion forests for heterogeneous treatment effects on survival
%% from randomised and real-world data" to the Biostatistics special
%% collection "Statistical Foundations of AI and Real-World Evidence
%% Generation" (submission deadline 31 July 2026).
%% Compile on its own; it is not part of main.tex or sm.tex.

\documentclass[11pt]{article}

\usepackage[a4paper, margin=1in]{geometry}
\usepackage{lmodern}
\usepackage{microtype}
\usepackage[colorlinks, urlcolor=blue]{hyperref}
\usepackage{xurl}

\setlength{\parindent}{0pt}
\setlength{\parskip}{0.7em}
\pagestyle{empty}

\begin{document}

\hfill \today

\medskip

The Editors and Guest Editors \\
Special Collection: Statistical Foundations of AI and Real-World Evidence
Generation \\
\textit{Biostatistics}

\medskip

Dear Editors,

We submit our manuscript ``Bayesian fusion forests for heterogeneous treatment
effects on survival from randomised and real-world data'' for consideration in
the special collection on Statistical Foundations of AI and Real-World Evidence
Generation.

Randomised trials identify treatment effects but are underpowered for subgroup
effects and often stop before the events of interest accumulate. Real-world
data are larger and follow patients for longer, but treatment there is not
randomised. Our manuscript develops the Bayesian fusion forest, a nonparametric
framework that combines the two sources without assuming the observational
source is unconfounded. A confounding function absorbs the bias from unmeasured
confounding, a shared baseline prognosis with a source-specific deviation
encodes the borrowing, and a hierarchical Dirichlet process mixture lets the
error distribution differ between the sources. Each component carries an
independent causal meaning and receives its own tree-ensemble prior. To our
knowledge the framework is the first that combines causal effect estimation,
data fusion and support for both right- and interval-censored outcomes.

The work speaks to several themes of the collection. It offers a statistical
perspective on combining experimental and real-world data, and it performs
causal inference with a machine-learning model for the purpose of real-world
evidence generation. Transportability across data sources is the assumption
that makes the fusion possible, and we state its cost explicitly. Uncertainty
quantification is intrinsic rather than added afterwards, and our cohort
records outcomes only at annual visits, so the event times are
interval-censored, one of the data imperfections the call names.

A simulation study shows efficiency gains over a trial-only analysis across
varying levels of confounding and between-source heterogeneity. We benchmark
against four machine-learning alternatives, and the Bayesian fusion forest
attains a
lower error than every alternative in every configuration while remaining the
only method whose intervals reach nominal coverage at a competitive width. We
apply the method to the ACTG~175 trial combined with the Multicenter AIDS
Cohort Study, where the fusion identifies a benefit of combination
antiretroviral therapy for nearly every patient while the trial alone is
inconclusive.

The manuscript is original, is not under consideration elsewhere, and all
authors have approved the submission. The authors declare no conflict of
interest. The software implementation, together with the code for the
simulation studies and the data analysis, is available at
\url{https://github.com/tijn-jacobs/FusionForests}.

The authors used generative artificial intelligence tools in support of
developing and debugging the accompanying code. All such output was reviewed and
verified by the authors, who take full responsibility for the content of this
work.

\medskip

Yours sincerely,

\medskip

Tijn Jacobs, on behalf of the authors \\
Department of Mathematics, Vrije Universiteit Amsterdam \\
\texttt{t.jacobs@vu.nl}

\end{document}
